% Generated draft: requires mathematical and source review.
\documentclass[11pt]{article}
\usepackage[margin=25mm]{geometry}
\usepackage{fontspec}
\setmainfont{Latin Modern Roman}
\usepackage{amsmath,amssymb,mathtools,unicode-math}
\setmathfont{Latin Modern Math}
\usepackage{enumitem,needspace,xurl,hyperref}
\setlength{\parindent}{0pt}
\setlength{\parskip}{5pt}
\setlength{\emergencystretch}{3em}
\newcommand{\N}{\mathbb{N}}
\newcommand{\Z}{\mathbb{Z}}
\newcommand{\Q}{\mathbb{Q}}
\newcommand{\R}{\mathbb{R}}
\newcommand{\C}{\mathbb{C}}
\newcommand{\F}{\mathbb{F}}
\newcommand{\A}{\mathbb{A}}
\newcommand{\PP}{\mathbb{P}}
\newcommand{\E}{\mathbb{E}}
\newcommand{\eps}{\varepsilon}
\newcommand{\dd}{\,\mathrm{d}}
\begin{document}
\section*{UnsolvedMath ID: 700530}
\textbf{Source problem code:} catalogue entry
\subsection{Definitions and mathematical statement}
Let \(K=\Q(\sqrt D)\) be the imaginary quadratic field intended in the source, where \(D<0\), and let \(\mathfrak o\) be its ring of integers. The source does not explicitly specify the domain of \(D\); no squarefreeness assumption is imposed, since its examples include \(D=-4\). Fix an embedding \(K\subset\C\), and write \(\overline{x}\) for complex conjugation. Interpret \(\pi\) as a prime element of \(\mathfrak o\) such that
\[
\pi\overline\pi=p
\]
for a positive rational prime \(p\). This element interpretation is suggested by the source's congruence notation, but is not explicitly distinguished there from prime-ideal terminology.

Define three-dimensional hyperbolic upper half-space by
\[
\mathfrak H^3=\{(z,r):z\in\C,\ r\in\R,\ r>0\}.
\]
For \(g=\begin{pmatrix}a&b\\c&d\end{pmatrix}\in\operatorname{SL}_2(\C)\), use the transformations supplied in the source:
\[
T_g(z,r)=\left(
\frac{(az+b)\overline{cz+d}+a\overline c\,r^2}
{|cz+d|^2+|c|^2r^2},
\frac{r}{|cz+d|^2+|c|^2r^2}
\right).
\]
Let
\[
\Gamma=\Gamma(K,\pi)=
\left\{
\begin{pmatrix}a&b\\c&d\end{pmatrix}\in\operatorname{SL}_2(\mathfrak o):
b\in\pi\mathfrak o
\right\},
\qquad
X=\Gamma\backslash\mathfrak H^3,
\]
where the quotient is by the orbits of these transformations. Let \([\Gamma,\Gamma]\) be the subgroup generated by all commutators, and put \(\Gamma^{\mathrm{ab}}=\Gamma/[\Gamma,\Gamma]\).

The source asserts that adjoining two vertices to \(X\) produces a compact three-dimensional space, denoted here by \(\overline X\). It asks to calculate the quantity presented as
\[
r(\pi)
=\dim_{\Q}H_1(\overline X;\Q)
=\dim_{\Q}\bigl(\Gamma^{\mathrm{ab}}\otimes_{\Z}\Q\bigr),
\]
where \(H_1\) denotes singular homology. Thus the unambiguously specified algebraic calculation is to determine
\[
r_{\mathrm{ab}}(K,\pi)
:=\dim_{\Q}\bigl(\Gamma(K,\pi)^{\mathrm{ab}}\otimes_{\Z}\Q\bigr)
\]
for every admissible field and prime element. The displayed equality with the homology of \(\overline X\) is retained as a supplied assertion, not established here. The source gives neither a topology for adjoining the two vertices nor additional hypotheses explaining the number of vertices or this homological identification. Consequently, the topological formulation and its equivalence to the algebraic formulation require clarification before the displayed common quantity can be treated as fully specified.

The source supplies the following examples: for \(D=-3\), the first nonzero value is said to occur for \(\pi\mid73\), with \(r(\pi)=1\); for \(D=-4\), it states \(r(\pi)=1\) for \(\pi\mid137\). These are preserved as source claims. The meaning of ``first'' is not explicitly defined.
\subsection{Short English statement}
For the indicated imaginary quadratic field and a prime element of rational-prime norm, determine the rational rank of the abelianization of the subgroup of \(\operatorname{SL}_2(\mathfrak o)\) whose upper-right entry is divisible by \(\pi\). The source also identifies this rank with the first rational Betti number of a compactification obtained by adding two vertices; the compactification and that identification require clarification. The supplied registry labels the problem open.
\subsection{Sources}
\begin{itemize}
\item Local UnsolvedMath registry, permanent ID \(700530\), source code \texttt{catalogue entry}, title \emph{Kourovka Notebook Problem 8.82}. The supplied record labels its status as open; this label is not an independently verified mathematical or literature assessment.
\item \emph{The Kourovka Notebook (Unsolved Problems in Group Theory)}, original problem 8.82, eighth issue (1982), proposed by H. Helling. The supplied background identifies the referenced edition as the twenty-first issue edition (2026), \texttt{arXiv:1401.0300v48}.
\item The supplied collection description attributes the Kourovka Notebook, Issues 1--21 (1965--2026), to editors E. I. Khukhro and V. D. Mazurov. The cited source text and its examples have not been independently inspected or verified for this draft.
\end{itemize}
\end{document}
